% ===========================================================================
% W11 Results section for the Overleaf paper (acmart, English).
% Paste after the Method section. Upload results/figures/*.png to Overleaf as figures/.
% All numbers are generated from results/ (runs.jsonl, bootstrap_pairs.json, tables/*.csv).
%
% CHECK BEFORE SUBMITTING
%  1. Section 2.3 of the Method says the TF-IDF vocabulary has 26,547 terms. The selected
%     configuration (E3-B) caps it at 20,000 features; state this in Method or Results.
%  2. Section 2.2 gives the *initial* ALS configuration (64 factors, lambda=0.1); the tuned
%     configuration (E2-07, lambda=0.01) is reported here.
%  3. Section 2.4 does not yet name the encoder (all-MiniLM-L6-v2, max length 256).
%  4. The abstract is still empty.
% ===========================================================================

\section{Experimental Setup}
\label{sec:setup}
All results are computed on the validation split of the official temporal warm-start
protocol: models are fitted on the training interactions only and ranked over the
22{,}753 warm candidate items for each of the 11{,}659 evaluable validation users, with
each user's training items removed. An item is relevant if the user rated it $\geq 4$ in the
warm validation data. The test set is not used. We report NDCG@10 (the primary metric and the
basis of model selection) and Recall@10 (tie-breaker), plus HitRate@10, MRR@10 and
Coverage@10 (the share of candidates that appear in at least one Top-10 list). Users are
grouped by the number of training interactions into low (1--4, $n=4{,}772$), medium (5--9,
$n=4{,}142$) and high ($\geq 10$, $n=2{,}745$) activity. The seed is 42.

Hyperparameters were tuned on validation only, following the search space fixed beforehand:
18 ALS configurations (factors $\{32,64,128\}$, $\lambda\in\{0.01,0.1,1\}$,
$\alpha\in\{10,40\}$), three TF-IDF configurations (5{,}000 or 20{,}000 features, with or
without bigrams) and three embedding configurations (maximum input length 128, 256 or 512 tokens),
so that the tuning budget of TF-IDF and the embedding model is equal. The configuration with the
highest NDCG@10 over all evaluable users is selected, with Recall@10 as tie-breaker. The
selected configurations are E2-07 (64 factors, $\lambda=0.01$, $\alpha=10$), E3-B (unigrams,
20{,}000 features) and E4-B (\texttt{all-MiniLM-L6-v2}, length 256). Significance is assessed
with a paired bootstrap over users (1{,}000 resamples, seed 42, 95\% CI).

\section{Results}
\label{sec:results}

\subsection{Comparison with the baselines}
Table~\ref{tab:main} and Figure~\ref{fig:models} compare the selected configurations at $K=10$.
ALS is the strongest model, followed by Popularity, TF-IDF and the embedding model. The order is
the same for NDCG, Recall, HitRate and MRR, and for $K=10, 20, 50$ (Table~\ref{tab:k},
Figure~\ref{fig:k}). ALS improves NDCG@10 over Popularity by 35\% and Recall@10 by 48\%;
TF-IDF is 18\% (NDCG) and 6\% (Recall) below Popularity, and the embedding model is a further
31\% (NDCG) and 32\% (Recall) below TF-IDF. All models are 20--46 times better than random
ranking (Recall@10 $\approx 10/22{,}753$).

\begin{table}[t]
\caption{Validation results at $K=10$ on all 11{,}659 evaluable users (selected configurations).}
\label{tab:main}
\small
\begin{tabular}{llccccc}
\toprule
Model & ID & Recall & NDCG & HitRate & MRR & Coverage \\
\midrule
Popularity & E1 & 0.0135 & 0.0092 & 0.0310 & 0.0130 & 0.0008 \\
ALS & E2-07 & 0.0200 & 0.0124 & 0.0396 & 0.0153 & 0.0421 \\
TF-IDF & E3-B & 0.0127 & 0.0075 & 0.0214 & 0.0082 & 0.7122 \\
Embedding & E4-B & 0.0087 & 0.0052 & 0.0149 & 0.0058 & 0.6221 \\
\bottomrule
\end{tabular}
\end{table}

\begin{figure}[t]
\centering
\includegraphics[width=\columnwidth]{figures/fig1_model_comparison.png}
\caption{NDCG@10 and Recall@10 of the selected configurations on the validation split.}
\Description{Bar charts of NDCG@10 and Recall@10 for Popularity, ALS, TF-IDF and the embedding model.}
\label{fig:models}
\end{figure}

\begin{table}[t]
\caption{Effect of $K$ on the selected models (R = Recall, N = NDCG).}
\label{tab:k}
\small
\begin{tabular}{lcccccc}
\toprule
Model & R@10 & N@10 & R@20 & N@20 & R@50 & N@50 \\
\midrule
Popularity & 0.0135 & 0.0092 & 0.0215 & 0.0115 & 0.0420 & 0.0162 \\
ALS & 0.0200 & 0.0124 & 0.0299 & 0.0154 & 0.0524 & 0.0206 \\
TF-IDF & 0.0127 & 0.0075 & 0.0195 & 0.0095 & 0.0348 & 0.0131 \\
Embedding & 0.0087 & 0.0052 & 0.0137 & 0.0067 & 0.0252 & 0.0093 \\
\bottomrule
\end{tabular}
\end{table}

\begin{figure}[t]
\centering
\includegraphics[width=\columnwidth]{figures/fig2_effect_of_k.png}
\caption{NDCG@K and Recall@K for $K\in\{10,20,50\}$.}
\Description{Line charts of NDCG@K and Recall@K for the four models; the model order is identical at every K.}
\label{fig:k}
\end{figure}

\subsection{Hyperparameter search}
Tuning changed the results very little (Figures~\ref{fig:grid} and~\ref{fig:lexsem}). The 18 ALS
configurations span NDCG@10 0.0116--0.0124; the selected E2-07 is only 0.2\% better than the initial
configuration, and E2-09 and E2-11 are practically tied with it (all 0.0124). Its regularization
$\lambda=0.01$ is the smallest value tried, so smaller values remain untested. Tuning improved
TF-IDF by 0.5\% (E3-B versus the initial E3) and left the embedding configuration unchanged (the
initial E4 equals E4-B). Every embedding configuration (NDCG@10 0.0047--0.0052) is below every TF-IDF
configuration (0.0071--0.0075), so the conclusion for RQ1 does not depend on the choice of
hyperparameters within this search space. The input length is not monotone: 256 tokens is best,
followed by 512 (0.0050) and 128 (0.0047).

\begin{figure}[t]
\centering
\includegraphics[width=\columnwidth]{figures/fig4_als_grid.png}
\caption{NDCG@10 of the 18 ALS configurations on the validation split.}
\Description{Heat map of NDCG@10 over factors and regularization/alpha combinations; values range from 0.0116 to 0.0124.}
\label{fig:grid}
\end{figure}

\begin{figure}[t]
\centering
\includegraphics[width=0.85\columnwidth]{figures/fig5_lexical_vs_semantic.png}
\caption{NDCG@10 of all TF-IDF (E3) and embedding (E4) configurations.}
\Description{Bar chart showing that all three embedding configurations score below all three TF-IDF configurations.}
\label{fig:lexsem}
\end{figure}

\subsection{RQ1: embedding versus TF-IDF}
Under the same item text, profile rule and candidate set, the sentence-embedding recommender does
\emph{not} achieve higher top-10 ranking quality than TF-IDF. The embedding model is lower on
NDCG@10 (0.0052 versus 0.0075) and Recall@10 (0.0087 versus 0.0127), at every $K$, and for
every configuration tried. Table~\ref{tab:boot} shows that the difference is statistically
significant: Embedding $-$ TF-IDF is $-0.0023$ NDCG@10 (95\% CI $[-0.0033, -0.0011]$) and $-0.0041$
Recall@10. The same table shows that ALS is significantly better than every other model and that the
embedding model is significantly below Popularity. TF-IDF is significantly below Popularity on NDCG@10
(upper confidence limit $-0.0001$) but not on Recall@10, where the interval includes zero.

\begin{table}[t]
\caption{Paired bootstrap of the difference in the mean metric at $K=10$ (all users; 95\% CI in brackets; $^{\mathrm{ns}}$ = interval includes zero).}
\label{tab:boot}
\small
\begin{tabular}{lcc}
\toprule
Comparison (A $-$ B) & $\Delta$NDCG@10 & $\Delta$Recall@10 \\
\midrule
ALS $-$ Popularity & +0.0032 [+0.0017, +0.0049] & +0.0065 [+0.0039, +0.0093] \\
ALS $-$ TF-IDF & +0.0049 [+0.0032, +0.0066] & +0.0073 [+0.0045, +0.0101] \\
ALS $-$ Embedding & +0.0072 [+0.0057, +0.0088] & +0.0113 [+0.0087, +0.0139] \\
TF-IDF $-$ Popularity & $-$0.0017 [$-$0.0034, $-$0.0001] & $-$0.0007 [$-$0.0033, +0.0018]$^{\mathrm{ns}}$ \\
Embedding $-$ Popularity & $-$0.0040 [$-$0.0055, $-$0.0025] & $-$0.0048 [$-$0.0070, $-$0.0025] \\
Embedding $-$ TF-IDF & $-$0.0023 [$-$0.0033, $-$0.0011] & $-$0.0041 [$-$0.0057, $-$0.0022] \\
\bottomrule
\end{tabular}
\end{table}

\subsection{RQ2: the gap across user activity}
Table~\ref{tab:group} and Figure~\ref{fig:group} break NDCG@10 down by training activity. The gap
between the semantic and lexical representation is largest for medium-activity users
($0.0039$), and smaller for low ($0.0012$) and high ($0.0017$) activity users. The bootstrap
(Table~\ref{tab:groupboot}) shows that the NDCG@10 gap is significant only for medium-activity
users; for low and high activity the interval includes zero, although the Recall@10 gap is
significant in all three groups. Both content-based models weaken for active users (TF-IDF
0.0089 $\to$ 0.0049 from medium to high; embedding 0.0066 $\to$ 0.0031 from low to high), whereas ALS
improves (0.0114 $\to$ 0.0148). A possible explanation is that the mean of many diverse item
vectors yields an over-general profile; we did not test this. For low-activity users ALS and
Popularity are indistinguishable (difference $+0.0003$, CI includes zero).

\begin{table}[t]
\caption{NDCG@10 by user activity group (low 1--4, medium 5--9, high $\geq 10$ training interactions).}
\label{tab:group}
\small
\begin{tabular}{lccc}
\toprule
Model & Low & Medium & High \\
\midrule
Popularity & 0.0111 & 0.0086 & 0.0070 \\
ALS & 0.0114 & 0.0121 & 0.0148 \\
TF-IDF & 0.0079 & 0.0089 & 0.0049 \\
Embedding & 0.0066 & 0.0050 & 0.0031 \\
TF-IDF $-$ Embedding & 0.0012 & 0.0039 & 0.0017 \\
\bottomrule
\end{tabular}
\end{table}

\begin{figure}[t]
\centering
\includegraphics[width=\columnwidth]{figures/fig3_activity_groups.png}
\caption{NDCG@10 by user activity group.}
\Description{Grouped bar chart of NDCG@10 for the four models in the low, medium and high activity groups.}
\label{fig:group}
\end{figure}

\begin{table}[t]
\caption{Paired bootstrap of the NDCG@10 difference by activity group (95\% CI; $^{\mathrm{ns}}$ = interval includes zero).}
\label{tab:groupboot}
\small
\begin{tabular}{llll}
\toprule
Comparison & Low & Medium & High \\
\midrule
Embedding $-$ TF-IDF & $-$0.0012 [$-$0.0028, +0.0002]$^{\mathrm{ns}}$ & $-$0.0039 [$-$0.0061, $-$0.0018] & $-$0.0018 [$-$0.0039, +0.0001]$^{\mathrm{ns}}$ \\
ALS $-$ Popularity & +0.0003 [$-$0.0022, +0.0026]$^{\mathrm{ns}}$ & +0.0036 [+0.0007, +0.0063] & +0.0078 [+0.0049, +0.0111] \\
TF-IDF $-$ Popularity & $-$0.0032 [$-$0.0059, $-$0.0005] & +0.0003 [$-$0.0025, +0.0034]$^{\mathrm{ns}}$ & $-$0.0021 [$-$0.0046, +0.0004]$^{\mathrm{ns}}$ \\
\bottomrule
\end{tabular}
\end{table}

\subsection{Further observations}
The two content-based models recommend very different items: the mean Jaccard@10 between the
TF-IDF and embedding lists is 0.151, 28.8\% of users share no
item, and only 0.14\% receive identical lists. The embedding model therefore does not
merely re-order the TF-IDF list; it selects other items that are less often correct. Accuracy and diversity
move in opposite directions: TF-IDF and the embedding model cover 71\% and 62\% of the catalogue,
against 4\% for ALS and 0.08\% for Popularity. Absolute values are small (Recall@10 of 1--2\%)
because each user has about two relevant items (54.5\% have exactly one) among 22{,}753
candidates, the data are sparse (density 0.0345\%) and the median user has five training interactions.

\section{Limitations}
\label{sec:limitations}
Results use one seed and one temporal split; variation across seeds is not measured. Only one encoder
(MiniLM) and one text field are compared, and user profiles are plain means of item vectors.
Significance is tested for the pairs in Table~\ref{tab:boot} and for the groups in
Table~\ref{tab:groupboot} without correction for multiple comparisons. Differences among the ALS
configurations are very small and were not tested. The test set has not been used; the final evaluation
of the selected configurations on it is future work.

\section{Reproducibility}
Code, configurations, results and figures are in the project repository
(\url{https://github.com/AndikaNadapdap/12S4054-G01-TrackE}). Every run records its configuration hash,
git commit and a clean-tree flag in \texttt{results/runs.jsonl}. Results are produced with
\texttt{scripts/run\_cp3\_validation.sh} ($K=10$), \texttt{scripts/run\_k\_sweep.sh} ($K=20,50$),
\texttt{scripts/run\_bootstrap.py} and \texttt{scripts/make\_figures.py}.
